\section{Introduction}
\begingroup
Interactive simulation, initial-condition studies, and geometry design often require solving similar flow problems many times. In airfoil design, engineers may repeatedly adjust shape parameters to identify the dominant flow modes associated with candidate geometries. In games, moving obstacles continually reshape the fluid domain and demand rapid visual feedback. Recomputing numerical bases and flow solutions for each configuration can make this interactive exploration costly. Reduced-order models (ROMs) reduce this cost by representing complex flows with a small set of spatial basis functions. A trained geometry-conditioned neural representation can query these modes directly, offering a route to near-real-time previews; once the reduced operators have been prepared, millisecond-scale coefficient updates support rapid flow evolution (\docref{Section}{sec:inference-cost}). The choice of basis determines what the reduced model can represent, how physical constraints are handled, and which computations can be reused. A useful fluid basis should therefore capture the main flow structures while supporting dynamics from different initial conditions.

For conventional spectral ROMs, however, the basis is usually computed by solving a Laplacian eigenvalue problem on a specific domain. Even a small change in the boundary or the position of an obstacle can change the eigenfunctions and require a new numerical solve. Analytical solutions are available for a few simple geometries, such as rectangles and circles, but are generally difficult to obtain for more complex shapes. This becomes a practical limitation during early-stage design, where geometry parameters may change repeatedly. Neural networks offer another approach. Instead of solving a separate eigenvalue problem for each geometry, we can learn a mapping from shape space to eigenspace, allowing basis functions to be evaluated at arbitrary spatial locations within a family of geometries. Our goal is to preserve the physical structure and accuracy of Laplacian Eigenfluids while reducing the need for repeated numerical solves. Our experiments show that learned bases can achieve accuracy comparable to numerical Laplacian bases, and sometimes perform better in selected cases. For repeated queries, the cost of training can also be shared across many evaluations.

Fluid model reduction has developed along two main directions: data-driven and operator-based methods. Snapshot-based approaches extract a low-dimensional space from existing simulations. \citet{treuille2006model} demonstrated interactive reduced fluid simulation, including interactions with moving obstacles. \citet{kim2013subspace} combined subspace integration with cubature to efficiently simulate variations of high-resolution flows. These methods use patterns in simulation data to reduce spatial computation and time-stepping costs. Spectral methods instead construct basis functions from differential operators. Laplacian Eigenfluids \citep{dewitt2012fluid} organizes flow structures by eigenvalues and provides a modal interpretation of viscous decay. Model-Reduced Variational Fluid Simulation \citep{liu2015model} further connects spectral reduction with geometric structure and harmonic fields. Unlike bases extracted from specific flow trajectories, operator-based modes are ordered by differential energy and have a clear physical interpretation. These properties provide a natural starting point for our neural representation.

Neural networks have introduced several other ways to reuse computations in fluid simulation. Deep Fluids learns a parameterized velocity generator and latent dynamics from simulation data \citep{kim2019deep}. DeepONet and Fourier Neural Operators learn mappings between function spaces \citep{lu2021deeponet,li2021fourier}, while Geo-FNO extends this approach to general geometries through learned deformations \citep{li2023geometry}. These methods mainly focus on learning how physical inputs map to flow solutions. We take a different approach: rather than learning the flow solution directly, we learn the spatial basis used to represent it. Once the basis and its reduced operators are prepared for a given geometry, they can be reused across different initial conditions by changing only the coefficient state. This also allows us to study spatial representation error and reduced dynamics error separately.

Coordinate-based neural fields provide continuous function values and spatial derivatives. Periodic activation functions help represent differentiable signals \citep{sitzmann2020implicit}, while Fourier features improve the learning of high-frequency variations \citep{tancik2020fourier}. Variational methods can then use physical energy to learn eigenfunctions. The Deep Ritz method minimizes energy functionals \citep{e2018deep}, Spectral Inference Networks learns multiple eigenfunctions through stochastic optimization \citep{pfau2019spectral}, and neural Rayleigh methods apply similar ideas to engineering eigenvalue problems \citep{rowan2025rayleigh}. More recently, Shape Space Spectra extended neural eigenfunction learning to parameterized geometries, using geometry-local sorting and causal orthogonalization to handle changes in spectral ordering \citep{chang2025shapespacespectra}. These works show that neural networks can learn bases defined by differential operators. For incompressible fluids, however, minimizing spectral energy alone is not enough. The basis must also satisfy incompressibility and wall impermeability, account for gauge freedom in the potential representation, and include the independent circulation fields introduced by domain topology.

Physical constraints can be built directly into neural representations. Neural Conservation Laws (NCL) constructs divergence-free fields through differential lifts of antisymmetric potentials \citep{richterpowell2022neural}, and other constrained neural basis methods use similar ideas for incompressible flows \citep{cheng2026decoupled}. Neural Kinematic Bases (NKB) learns velocity basis functions directly, using losses for smoothness, divergence, wall alignment, and orthogonality to support mesh-free flow fitting and animation \citep{liu2025neural}. These studies show that learning fluid bases with neural networks is feasible. However, soft penalties do not strictly enforce all physical constraints, and they do not directly organize the learned functions according to the spectrum of a differential operator.

We aim to reduce how much of the physics the network must learn from penalties. Instead of relying on the loss function to enforce incompressibility, we build it into the representation and use energy minimization to learn the spectral modes. In two dimensions, a scalar streamfunction generates a divergence-free velocity field; in three dimensions, we use a vector potential. Potential boundary conditions enforce wall impermeability, geometry-local orthogonalization separates the modes, and Rayleigh minimization determines their spectral ordering. The training does not require flow snapshots, since the basis is defined by the geometry and differential operator. The neural network provides continuous spatial functions, the physical construction restricts the space of possible solutions, and the energy objective selects the low-frequency modes. We also compare our approach with NKB under the same training budget and Galerkin evolution framework to examine how these structural choices affect accuracy.

Learning in potential space offers another advantage: it can reduce the number of spatial derivatives needed during training. However, this requires a consistent treatment of gauge freedom. For example, adding a constant to a periodic streamfunction in two dimensions does not change the velocity. In three dimensions, certain curl-free components can also be added to a vector potential without changing its curl. Although these changes leave the physical velocity unchanged, they affect the potential mass norm and can therefore distort the Rayleigh objective. We address this by selecting a suitable gauge representative before normalization and optimization. Using the classical compatible positive-spectrum correspondence \citep{arnold2006finite}, we connect the potential-space and velocity-space objectives and derive a spectral-variance relation that explains why they assign different weights to high-frequency components in approximate modes.

Our experiments also reveal a practical difference between two- and three-dimensional learning. In 2D, the potential objective requires fewer spatial derivatives and reduces training time while maintaining good accuracy. In 3D, however, selecting a consistent vector-potential gauge introduces a substantial computational cost, while the accuracy gains over numerical Laplacian bases are relatively small in our tests. We therefore view potential-space learning as an alternative with a clear mathematical basis, rather than an objective that is always preferable. Its usefulness depends on the problem and the cost of gauge selection. We further examine different auxiliary solvers to understand how gauge accuracy depends on spatial resolution and operator consistency.

Beyond the positive spectral modes, flow in complex domains contains independent circulation fields determined by topology. Hodge theory shows that harmonic velocity fields can carry circulation around obstacles without introducing vorticity \citep{arnold2006finite}. This structure has been studied in spectral variational ROMs \citep{liu2015model}, Fluid Cohomology \citep{yin2023fluid}, and Hodge-based neural operator learning \citep{zheng2026hodge}. In our framework, we learn these harmonic fields through fixed boundary traces and Dirichlet energy minimization, then convert them into velocity fields through differential lifts. Together, the positive modes and harmonic coordinates represent both vortical motion and independent circulation around obstacles.

Circulation is especially important in flow around lifting airfoils. In idealized airfoil flow, it is closely related to the Kutta condition and the generation of lift. A basis containing only positive spectral modes cannot represent the independent circulation needed for such a flow, regardless of how accurately those modes are learned. Harmonic completion therefore provides more than a mathematical correction; it adds degrees of freedom needed to describe certain physical phenomena. However, when the total number of modes is fixed, adding harmonic fields also means removing the same number of positive modes. Although this makes the representation topologically complete, it does not always improve energy capture. The benefit depends on how much of the target flow lies in the vortical and harmonic sectors.

\begin{samepage}
Our main contributions are as follows:
\begin{enumerate}[leftmargin=*,itemsep=3pt,topsep=4pt]
\item \textbf{Physically constrained neural eigenfluid bases.} We introduce a continuous, geometry-conditioned spectral representation that combines differential lifts, impermeable boundary conditions, and geometry-local orthogonalization. The resulting basis satisfies key kinematic constraints by construction. For a selected geometry, reusable Galerkin operators support autonomous coefficient evolution and flow reconstruction at arbitrary spatial locations.
\item \textbf{Potential-space spectral learning with gauge selection.} We formulate a Rayleigh learning objective on selected potential representatives and derive a spectral-variance relation connecting potential-space and velocity-space objectives through their shared positive spectrum. We also examine the accuracy and computational cost of both approaches in two and three dimensions.
\item \textbf{Neural harmonic circulation fields.} We learn independent harmonic coordinates using fixed boundary traces and Dirichlet energy minimization, completing the flow representation in multiply connected domains. Through fixed-capacity comparisons and controlled tests of boundary-lift conditioning, we study how these coordinates affect energy capture and flow evolution.
\end{enumerate}
\end{samepage}

\par\endgroup
